\chapter{Brushing Your Teeth}
\index{Brushing Your Teeth}

\section*{Definition and Overview}
Brushing the teeth is the most important daily habit for maintaining oral health. It removes the sticky film of plaque -- a community of bacteria mixed with food debris -- that constantly forms on tooth surfaces. If plaque is not removed regularly it hardens into tartar, irritates the gums, and drives tooth decay and gum disease. Correct brushing, performed twice a day with a fluoride toothpaste, is the cornerstone of preventing the two most common dental conditions: dental caries (cavities) and periodontal (gum) disease. Effective brushing also improves breath and can be complemented by interdental cleaning, which reaches the spaces a brush cannot.

\section*{Epidemiology}
Tooth decay and gum disease are among the most widespread health conditions worldwide. Tooth decay affects the majority of adults at some point in life and is very common in children, although rates have fallen in countries with water fluoridation. Some degree of gum disease affects roughly half of adults, with the severe destructive form (periodontitis) present in a smaller but significant proportion, concentrated in middle and old age. People with poor brushing habits, frequent sugar intake, smoking, or limited access to dental care have substantially higher rates of both diseases, and oral disease remains one of the most common reasons for preventable health-care attendance.

\section*{Aetiology and Pathophysiology}
Plaque bacteria metabolise sugars from the diet and produce acid that demineralises tooth enamel, causing decay. The same bacteria irritate the gingival tissues; when plaque is not removed, the gums become inflamed (gingivitis), and with time the deeper supporting structures -- the periodontal ligament and alveolar bone -- are destroyed (periodontitis).

Poor brushing contributes to disease in several distinct ways.

\begin{itemize}
    \item \textbf{Insufficient frequency or duration:} infrequent or rushed brushing leaves plaque in place for longer, allowing more acid damage
    \item \textbf{Missed surfaces:} failure to brush between the teeth and along the gum line allows decay to start in exactly those areas
    \item \textbf{Excessive force:} brushing too hard, or with a worn or stiff brush, damages gum tissue and wears away enamel
    \item \textbf{No fluoride:} not using a fluoride toothpaste removes the protective, remineralising effect that helps repair early damage
    \item \textbf{Evening neglect:} skipping brushing before bed permits acid attack to continue unchecked through the night
\end{itemize}

\section*{Clinical Features}
Problems that develop from inadequate brushing are usually noticed as the following signs.

\begin{itemize}
    \item Bleeding gums during brushing or flossing
    \item Red, swollen, or tender gums
    \item Persistent bad breath (halitosis)
    \item Tooth sensitivity to hot, cold, or sweet foods
    \item Visible holes, dark spots, or rough areas on the teeth
    \item Receding gums or loose, drifting teeth in advanced gum disease
    \item Pain or discomfort when biting
\end{itemize}

\section*{Differential Diagnosis}
The signs of poor oral hygiene overlap with other conditions, so professional assessment is important.

\begin{itemize}
    \item \textbf{Tooth decay versus staining or cracks:} enamel discoloration and craze lines can be mistaken for cavities
    \item \textbf{Gingivitis versus early periodontitis:} reversible gum inflammation versus disease affecting the supporting bone
    \item \textbf{Hormonal gum changes:} pregnancy-related gingivitis and bleeding
    \item \textbf{Bleeding disorders:} gum bleeding from thrombocytopenia or clotting defects
    \item \textbf{Oral mucosal disease:} mouth ulcers, oral thrush, or lichen planus
    \item \textbf{Referred pain:} sinus pain or tooth grinding mistaken for toothache
    \item \textbf{Oral cancer:} a non-healing sore, lump, or persistent ulcer must be assessed promptly
\end{itemize}

\section*{When to Seek Medical Care}
See a dentist if bleeding gums do not settle with improved brushing, if bad breath persists, if there is toothache or sensitivity, or if a sore in the mouth fails to heal within two to three weeks. Regular check-ups every six to twelve months allow problems to be detected and treated early.

\textbf{Contact your local emergency services for:}
\begin{itemize}
    \item Swelling of the face or jaw that makes it hard to breathe or swallow
    \item Severe dental pain with rapidly spreading swelling and fever
    \item Heavy bleeding from the mouth that will not stop
    \item Any dental problem accompanied by signs of serious infection, such as high fever or confusion
\end{itemize}

\section*{Diagnosis and Investigations}
A dentist diagnoses most problems by clinical examination, with investigations used where needed.

\begin{itemize}
    \item \textbf{Clinical examination:} inspection of the teeth and gums, with a probe to detect decay and measure the depth of the gum pockets
    \item \textbf{Dental X-rays:} to reveal decay between the teeth, bone loss around the roots, and problems below the gum line
    \item \textbf{Disclosing tablets:} a dye that stains plaque, showing where brushing has been missed
    \item \textbf{Periodontal probing:} recording of pocket depths and bleeding points to stage gum disease
    \item \textbf{Review of technique:} observation of brushing by the dentist or hygienist, with correction of errors
\end{itemize}

\section*{Management}
The main treatment for plaque-related disease is improving daily oral hygiene, supported by professional care.

\begin{itemize}
    \item \textbf{Brushing:} brush twice a day for two minutes with a fluoride toothpaste and a soft-bristled brush
    \item \textbf{Technique:} gentle circular strokes directed at the gum line, covering the outer, inner, and biting surfaces of every tooth
    \item \textbf{Interdental cleaning:} floss or interdental brushes once a day, because a toothbrush cannot reach between the teeth
    \item \textbf{Professional cleaning:} scaling and polishing to remove tartar that brushing cannot
    \item \textbf{Fluoride therapy:} additional fluoride applications or high-strength toothpaste for people at high risk of decay
    \item \textbf{Restorations:} fillings or crowns to repair cavities once they have formed
    \item \textbf{Periodontal treatment:} root surface debridement, and occasionally surgery, for established periodontitis
    \item \textbf{Brush replacement:} every three months, or after an illness
\end{itemize}

\section*{Complications}
Poor oral hygiene, if uncorrected, leads to progressive and sometimes serious problems.

\begin{itemize}
    \item Decay extending into the pulp, causing severe pain, abscess, and possible need for root-canal treatment or extraction
    \item Gingivitis progressing to periodontitis, with gum recession, bone loss, and drifting teeth
    \item Loose teeth and tooth loss
    \item Dental abscess with facial swelling and, rarely, spread of infection into deep facial spaces or the bloodstream
    \item Damage from over-brushing, including receded gums and worn enamel
    \item Associations between periodontal disease and systemic conditions such as diabetes and cardiovascular disease
    \item Social and psychological effects of bad breath and damaged teeth on confidence
\end{itemize}

\section*{Prognosis and Outcomes}
With regular, correct brushing, most people can prevent tooth decay and gum disease entirely. Early gingivitis is fully reversible with improved cleaning. Once periodontitis has caused bone loss the damage cannot be reversed, but the disease can usually be halted and the teeth retained for life. Cavities do not heal once formed, but small lesions can be caught early and restored simply. The earlier problems are found, the simpler and less costly the treatment, which is why routine dental attendance is so valuable.

\section*{Prevention and Self-Care}
Good prevention rests on daily habits and regular professional care.

\begin{itemize}
    \item Brush twice a day for two minutes with a fluoride toothpaste
    \item Clean between the teeth daily with floss or interdental brushes
    \item Limit sugary snacks and drinks, and avoid brushing immediately after acidic foods
    \item Do not smoke, and limit alcohol intake
    \item Drink fluoridated water where available
    \item Replace the toothbrush every three months
    \item See a dentist and hygienist regularly for check-ups and professional cleaning
    \item Teach and supervise children's brushing until they are around eight years old
\end{itemize}

\section*{Sources and Further Reading}
The following organisations provide reliable information on tooth brushing, oral hygiene, and dental health.

\begin{itemize}
    \item NHS. \textit{How to keep your teeth clean and other oral hygiene advice.} https://www.nhs.uk/live-well/healthy-teeth-and-gums/
    \item Centers for Disease Control and Prevention. \textit{Oral health: brushing, flossing, and prevention information.} https://www.cdc.gov/oral-health/
    \item World Health Organization. \textit{Oral health fact sheets and global guidance.} https://www.who.int/
    \item MedlinePlus. \textit{Teeth and gum care information.} https://medlineplus.gov/teethandgumcare.html
    \item Mayo Clinic. \textit{Oral health and tooth brushing advice.} https://www.mayoclinic.org/
    \item Australian Dental Association. \textit{Oral health tips and tooth brushing information.} https://www.teeth.org.au/
\end{itemize}
